% answers_short.tex —— 2010 年·普通高等学校招生考试·山东卷·文科 · 简短答案速查
% 由 tools/make_short_answers.py 自 answers.tex 的 \gksolhead/\gkanswer 提取，
% 并与 2010/2010_数学真题_山东文_答案要点.md 逐题比对后生成，勿手改（改源 MD 后重跑）。

\begin{gknotice}
\textbf{使用说明}：本节只给最终答案便于快速核对，完整推导见解析版。答案为编写者独立推导，\textbf{非官方标准答案}。
\end{gknotice}

\gkpart{选择题}
\begin{center}\begin{tabular}{ccccccccccccc}
\toprule
题号 & 1 & 2 & 3 & 4 & 5 & 6 & 7 & 8 & 9 & 10 & 11 & 12 \\
\midrule
答案 & C & B & A & D & A & B & C & C & B & D & A & B \\
\bottomrule
\end{tabular}\end{center}

\gkpart{填空题}
\begin{itemize}[leftmargin=2.6em]
  \item[\textbf{13.}] $-\dfrac54$。
  \item[\textbf{14.}] $3$。
  \item[\textbf{15.}] $\dfrac\pi6$。
  \item[\textbf{16.}] $x+y-3=0$。
\end{itemize}

\gkpart{解答题}
\begin{itemize}[leftmargin=2.6em]
  \item[\textbf{17.}] （1）$\omega=1$；（2）最小值 $1$。
  \item[\textbf{18.}] （1）$a_n=2n+1$，$S_n=n^2+2n$；（2）$T_n=\dfrac{n}{4(n+1)}$。
  \item[\textbf{19.}] （1）$\dfrac13$；（2）$\dfrac{13}{16}$。
  \item[\textbf{20.}] （1）证明；（2）$\dfrac14$。
  \item[\textbf{21.}] （1）$y=x+\ln2$；（2）$a<0$ 或 $a=0$ 时，减区间 $(0,1)$、增区间 $(1,+\infty)$；$0<a<\dfrac12$ 时，减区间 $(0,1)$ 与 $\left(\dfrac1a-1,+\infty\right)$，增区间 $\left(1,\dfrac1a-1\right)$；$a=\dfrac12$ 时在 $(0,+\infty)$ 上单调递减。
  \item[\textbf{22.}] （1）$\dfrac{x^2}{2}+y^2=1$；（2）① 证明，② $P(0,2)$ 或 $P\left(\dfrac54,\dfrac34\right)$。
\end{itemize}
